\documentclass[11pt,a4paper]{article}
\usepackage[utf8]{inputenc}
\usepackage[T1]{fontenc}
\usepackage[french]{babel}
\usepackage[margin=2cm]{geometry}
\usepackage{xcolor}
\usepackage{enumitem}
\usepackage{titlesec}
\usepackage{setspace}
\usepackage{parskip}

\definecolor{accent}{RGB}{234, 88, 12}
\definecolor{timing}{RGB}{37, 99, 235}
\definecolor{tip}{RGB}{22, 163, 74}

\titleformat{\section}{\Large\bfseries\color{accent}}{}{0pt}{}[\vspace{-4pt}\color{accent}\rule{\textwidth}{1.2pt}]
\titleformat{\subsection}{\normalsize\bfseries\color{timing}}{}{0pt}{}

\setstretch{1.25}

\begin{document}

\begin{center}
{\Huge\bfseries\color{accent} Notes de Présentation Intégrales --- PFE}\\[8pt]
{\Large Chaque slide (1 à 39) détaillé}\\[6pt]
{\large Salah Eddine AHADAF --- Juin 2026}\\[6pt]
{\color{timing}\textbf{Durée cible : $\sim$16--17 minutes}}
\end{center}

\vspace{0.5cm}

%=============================================================================
\section{Slide 1 --- Page de Garde}
\subsection{Durée : $\sim$10 secondes}
%=============================================================================
Bonjour membres du jury, bonjour à tous. Je vous présente aujourd'hui mon projet de fin d'études portant sur le développement d'une plateforme d'intelligence territoriale, visant à optimiser le réseau d'agences d'Attijariwafa Bank grâce à la donnée.

%=============================================================================
\section{Slide 2 --- Plan de la Présentation}
\subsection{Durée : $\sim$15 secondes}
%=============================================================================
Voici le plan de la présentation. Nous commencerons par l'introduction du contexte. Ensuite, je vous présenterai l'implémentation de la solution et ses trois scénarios analytiques principaux, avant de passer à une démonstration du tableau de bord. Enfin, nous terminerons par la conclusion et les perspectives.

%=============================================================================
\section{Slide 3 --- Intercalaire : Introduction (Contexte et Organisme)}
\subsection{Durée : $\sim$5 secondes}
%=============================================================================
Commençons tout de suite par le contexte du projet et l'organisme d'accueil.

%=============================================================================
\section{Slide 4 --- Attijariwafa Bank}
\subsection{Durée : $\sim$25 secondes}
%=============================================================================
J'ai effectué mon stage au sein d'Attijariwafa Bank. Premier groupe bancaire au Maghreb et en Afrique de l'Ouest, la banque dispose d'un réseau physique massif avec plus de 20 000 collaborateurs. Ce qui nous intéresse particulièrement ici, ce sont ses plus de 3 800 agences réparties sur le territoire marocain.

%=============================================================================
\section{Slide 5 --- Contexte du Projet}
\subsection{Durée : $\sim$35 secondes}
%=============================================================================
Le secteur bancaire évolue rapidement. Bien que la digitalisation modifie la donne, la présence physique reste cruciale. Cependant, la dynamique urbaine au Maroc change très vite : nouveaux quartiers périphériques, étalement urbain, centres-villes qui se saturent. Les méthodes traditionnelles d'implantation, souvent basées sur l'intuition, montrent leurs limites. On observe des agences qui se cannibalisent et des zones à fort potentiel ignorées. L'opportunité est donc de tirer parti de la donnée géospatiale pour rationaliser l'évaluation du réseau.

%=============================================================================
\section{Slide 6 --- Intercalaire : Problématique et Objectifs}
\subsection{Durée : $\sim$5 secondes}
%=============================================================================
Cela nous amène naturellement à la problématique.

%=============================================================================
\section{Slide 7 --- Problématique}
\subsection{Durée : $\sim$30 secondes}
%=============================================================================
Comment aider Attijariwafa Bank à analyser et optimiser son réseau d'agences à l'aide des données géospatiales ? Le but est de passer d'une gestion intuitive à une gestion mathématiquement justifiable auprès de la direction. Cette problématique se divise en trois axes : évaluer l'existant de manière neutre, optimiser les performances par des fermetures ou des fusions, et recommander intelligemment de nouvelles implantations.

%=============================================================================
\section{Slide 8 --- Objectifs du Projet}
\subsection{Durée : $\sim$30 secondes}
%=============================================================================
L'objectif est de créer une véritable plateforme d'Intelligence Territoriale. Pour y arriver, nous devons collecter les points d'intérêt depuis plusieurs sources, calculer des scores normalisés, segmenter le réseau par un algorithme de clustering, détecter la cannibalisation, utiliser le Machine Learning pour recommander des sites, et classer les agences sur une matrice BCG spatiale.

%=============================================================================
\section{Slide 9 --- Intercalaire : Implémentation (Conception Fonctionnelle)}
\subsection{Durée : $\sim$5 secondes}
%=============================================================================
Voyons maintenant comment nous avons conçu cette solution sur le plan fonctionnel.

%=============================================================================
\section{Slide 10 --- Exigences Fonctionnelles}
\subsection{Durée : $\sim$25 secondes}
%=============================================================================
Nous avons organisé les exigences autour de trois pôles. L'axe "Évaluer", qui gère le scraping et le calcul des sept scores. L'axe "Optimiser", qui repose sur un modèle prédictif XGBoost pour catégoriser les agences sur la matrice stratégique. Et l'axe "Recommander", qui s'occupe de la sélection spatiale de nouveaux sites avec filtrage et contraintes.

%=============================================================================
\section{Slide 11 --- Workflow Global}
\subsection{Durée : $\sim$30 secondes}
%=============================================================================
Ce diagramme illustre le flux de bout en bout. On part de l'ingestion des données brutes, avec une phase critique de dédoublonnage géographique — car des points d'intérêt doublonnés fausseraient totalement l'analyse stratégique finale. On passe ensuite par le scoring et le clustering, puis on alimente le moteur de sélection de sites et le module d'optimisation BCG.

%=============================================================================
\section{Slide 12 --- Intercalaire : Implémentation (Conception Technique)}
\subsection{Durée : $\sim$5 secondes}
%=============================================================================
Côté conception technique, l'architecture repose sur un pipeline robuste.

%=============================================================================
\section{Slide 13 --- Pipeline de Données et Features}
\subsection{Durée : $\sim$55 secondes}
%=============================================================================
Nous avons agrégé quatre sources de données : OpenStreetMap, Overture, Google Maps et les données internes AWB. L'extraction des données de Google Maps a d'ailleurs nécessité le développement d'un scraper personnalisé avec Selenium, intégrant un système de points de sauvegarde (checkpoints) pour résister aux coupures réseaux lors du traitement massif de la grille marocaine.

Suivant la première loi de la géographie de Tobler, qui stipule que "ce qui est proche a plus d'impact que ce qui est loin", nous utilisons une agrégation radiale autour de l'agence. L'attractivité commerciale, résidentielle, ou l'accessibilité viennent augmenter le potentiel de l'agence. À l'inverse, les scores de concurrence et de cannibalisation sont des pénalités soustraites lors du calcul du score composite global.

%=============================================================================
\section{Slide 14 --- Outils et Technologies}
\subsection{Durée : $\sim$25 secondes}
%=============================================================================
L'ensemble est codé en Python. Streamlit sert de framework frontend, la cartographie repose sur Folium, et la visualisation sur Plotly. Côté Machine Learning, j'ai utilisé Scikit-learn, XGBoost, et la structure algorithmique cKDTree de SciPy, absolument vitale pour réaliser nos requêtes de distance spatiale à grande échelle en temps optimal $O(N \log N)$.

%=============================================================================
\section{Slide 15 --- Intercalaire : Les principaux Scénarios}
\subsection{Durée : $\sim$5 secondes}
%=============================================================================
À partir de cette base technique, le cœur de la solution déploie trois grands scénarios d'usage.

%=============================================================================
\section{Slide 16 --- Intercalaire : SC1 Évaluation du Réseau}
\subsection{Durée : $\sim$5 secondes}
%=============================================================================
Le premier scénario est l'évaluation approfondie du réseau existant.

%=============================================================================
\section{Slide 17 --- Analyse Exploratoire et PCA}
\subsection{Durée : $\sim$40 secondes}
%=============================================================================
En analysant la distribution du score composite, on observe que le réseau est très hétérogène. Pour structurer cette donnée, j'ai appliqué un algorithme K-Means. Le choix de 5 clusters s'est imposé comme étant l'optimum mathématique et métier. La projection PCA montre ces cinq clusters : l'axe principal capture 72,5\% de la variance, séparant la densité urbaine des zones rurales, et l'axe secondaire différencie la vocation du quartier, opposant l'activité commerciale à l'activité résidentielle.

%=============================================================================
\section{Slide 18 --- Les 5 Profils d'Agences}
\subsection{Durée : $\sim$35 secondes}
%=============================================================================
Ces 5 profils ont une vraie signification métier. Les "High Performers" représentent un tiers du réseau. Les "Commercial Hubs" s'opposent aux "Residential Anchors" : l'un est en zone d'affaires générant des flux diurnes, l'autre en zone de vie nécessitant une relation client de proximité. On a aussi 17\% d'agences "At Risk", avec une très forte cannibalisation interne, et enfin 10\% d'"Underperformers", fortement isolés, comme on le voit sur le graphique radar rétracté.

%=============================================================================
\section{Slide 19 --- Benchmarking : Modèles de Régression}
\subsection{Durée : $\sim$40 secondes}
%=============================================================================
Pour prédire le revenu généré par une agence en fonction de son environnement, j'ai comparé plusieurs algorithmes. Le Random Forest l'emporte avec un R-carré de 0,884, surpassant largement les modèles linéaires. Cela s'explique techniquement : la relation entre le nombre de commerces et les revenus bancaires n'est pas linéaire. Il y a des effets de seuil et de saturation que seuls les modèles arborescents, nativement plus robustes face aux valeurs extrêmes, parviennent à capturer.

%=============================================================================
\section{Slide 20 --- Importance des Caractéristiques}
\subsection{Durée : $\sim$25 secondes}
%=============================================================================
En observant l'importance des variables calculée par XGBoost, une forte intuition métier se confirme : le score commercial est de très loin le prédicteur dominant du revenu. Une agence implantée près des malls et des zones de chalandise génère beaucoup plus de flux de prospects spontanés, impactant directement les ventes.

%=============================================================================
\section{Slide 21 --- Benchmarking : Classification}
\subsection{Durée : $\sim$40 secondes}
%=============================================================================
J'ai également modélisé le problème en classification avec 3 catégories : KEEP, REVIEW ou CLOSE. Ici, c'est la régression logistique qui performe le mieux avec 83\% de précision, battant même XGBoost. C'est cohérent : les classes extrêmes "KEEP" et "CLOSE" sont très linéairement séparables. Le modèle n'hésite que sur la classe intermédiaire "REVIEW", qui par nature représente une zone de transition financière difficile à délimiter strictement.

%=============================================================================
\section{Slide 22 --- Intercalaire : SC2 Sélection Nouveaux Sites}
\subsection{Durée : $\sim$5 secondes}
%=============================================================================
Passons au deuxième scénario : le moteur de sélection de nouveaux sites.

%=============================================================================
\section{Slide 23 --- Pipeline de Sélection}
\subsection{Durée : $\sim$50 secondes}
%=============================================================================
Ce moteur génère d'abord une grille de points tous les 500m sur le pays, un espacement choisi pour équilibrer la précision géographique et le temps de calcul. On calcule les scores pour ces candidats, puis on mesure leur similarité cosinus avec nos meilleures agences actuelles. L'avantage du cosinus, c'est qu'il compare l'angle des vecteurs, donc le "profil" de l'environnement, sans être biaisé par le volume total. Le Random Forest prédit le revenu, et l'algorithme sélectionne les 20 meilleurs sites, avec une contrainte métier stricte : 2 sites maximum par ville et au moins 700m d'écart.

%=============================================================================
\section{Slide 24 --- Résultats Sélection de Sites}
\subsection{Durée : $\sim$20 secondes}
%=============================================================================
Sur la plateforme, cela se traduit par ces étoiles orange pointant les 20 emplacements optimaux. Au survol, on accède immédiatement au revenu potentiel prédit et au score de similarité. Le directeur réseau dispose ainsi d'une justification mathématique solide pour déclencher une étude de marché locale.

%=============================================================================
\section{Slide 25 --- Intercalaire : SC3 Optimisation du Réseau}
\subsection{Durée : $\sim$5 secondes}
%=============================================================================
Le troisième scénario concerne l'optimisation stratégique et l'analyse de la valeur.

%=============================================================================
\section{Slide 26 --- Gap Analysis Concept}
\subsection{Durée : $\sim$50 secondes}
%=============================================================================
C'est le concept le plus actionnable du projet : la Gap Analysis. Son principe est de totalement découpler la géographie du management. 
Le modèle XGBoost, aveugle aux équipes, prédit ce que l'emplacement brut "devrait" rapporter. En divisant le revenu réel par cet attendu, on isole un ratio de performance de l'équipe. En croisant ce ratio managérial et le potentiel de la zone, on place chaque agence sur notre matrice BCG : Flagship, Review, Monitor ou Close. Une cinquième action, Merge, gère les consolidations de réseau.

%=============================================================================
\section{Slide 27 --- Matrice BCG Résultats}
\subsection{Durée : $\sim$20 secondes}
%=============================================================================
La matrice spatiale se lit ainsi : l'axe X représente le potentiel estimé par le modèle, l'axe Y le revenu réel. La diagonale parfaite représente les agences dont le revenu correspond exactement à la qualité de l'emplacement. Tout ce qui est au-dessus est en sur-performance d'équipe, tout ce qui est en-dessous est en sous-performance.

%=============================================================================
\section{Slide 28 --- Répartition Stratégique}
\subsection{Durée : $\sim$45 secondes}
%=============================================================================
Au niveau national, l'opportunité majeure réside dans les 86 agences classées "Review". Ce sont des emplacements en or qui génèrent de mauvais revenus. Cela signale un problème purement opérationnel. Plutôt que d'ouvrir une nouvelle agence, la banque pourrait simplement rénover ces agences "Review", y installer de nouveaux GAB, ou restructurer l'équipe commerciale pour les transformer en "Flagship". À l'inverse, les 513 agences "Close", mal placées et non-rentables, sont des candidates à la rationalisation.

%=============================================================================
\section{Slide 29 --- Fusions Géospatiales}
\subsection{Durée : $\sim$40 secondes}
%=============================================================================
L'algorithme va plus loin avec son arbre cKDTree. Pour chaque agence à fermer, il vérifie si un "Flagship" se trouve à moins de 1,2 km. Si c'est le cas, la recommandation devient une "Fusion". D'un point de vue rentabilité, cela permet de résilier un bail et de réduire drastiquement les dépenses d'exploitation (OpEx), tout en rapatriant le fonds de commerce vers l'agence modèle à proximité, maintenant ainsi la part de marché intacte.

%=============================================================================
\section{Slide 30 --- Intercalaire : Démonstration}
\subsection{Durée : $\sim$5 secondes}
%=============================================================================
Voyons rapidement comment cela se présente dans le tableau de bord livré.

%=============================================================================
\section{Slide 31 --- Architecture Platforme}
\subsection{Durée : $\sim$15 secondes}
%=============================================================================
Le produit fini est composé de 10 onglets couvrant tout : cartographie, classements, détections de fusions, audit individuel d'agence, matrice BCG, jusqu'à l'outil de scraping intégré.

%=============================================================================
\section{Slide 32 --- Module Map}
\subsection{Durée : $\sim$15 secondes}
%=============================================================================
Voici la vue cartographique avec nos 1 198 agences. Les décideurs disposent de filtres dynamiques, et le code couleur permet d'identifier d'un seul coup d'œil les agences saines en vert, ou celles nécessitant une action prioritaire en rouge ou orange.

%=============================================================================
\section{Slide 33 --- Rankings et Clusters}
\subsection{Durée : $\sim$15 secondes}
%=============================================================================
Ces modules offrent une lecture tabulaire. La direction peut y filtrer les agences d'une ville, isoler le Top 20 national, ou explorer la distribution des profils Machine Learning pour mieux comprendre la dynamique commerciale locale.

%=============================================================================
\section{Slide 34 --- Mergers et Cities}
\subsection{Durée : $\sim$15 secondes}
%=============================================================================
Le module Mergers trace des liaisons rouges entre les agences en situation de cannibalisation. Le module Cities génère des radars comparatifs pour opposer instantanément la compétitivité entre deux villes, facilitant l'allocation budgétaire régionale.

%=============================================================================
\section{Slide 35 --- Deep Dive et Data Enrichment}
\subsection{Durée : $\sim$30 secondes}
%=============================================================================
L'onglet Deep Dive est une fiche d'audit complète d'une agence cliquée. Enfin, l'onglet Data Enrichment permet aux équipes métier de relancer un scraping Google Maps. Techniquement, ce processus tourne en tâche de fond asynchrone, ce qui permet à l'interface Streamlit de rester fluide et réactive pendant que la base de données se met à jour.

%=============================================================================
\section{Slide 36 --- Intercalaire : Conclusion et Perspectives}
\subsection{Durée : $\sim$5 secondes}
%=============================================================================
Pour conclure cette présentation.

%=============================================================================
\section{Slide 37 --- Conclusion}
\subsection{Durée : $\sim$40 secondes}
%=============================================================================
Ce projet a abouti au déploiement réussi d'une plateforme d'Intelligence Territoriale complète pour Attijariwafa Bank. 

Nous avons remplacé l'intuition par un scoring spatial objectif. Nous avons fourni un moteur algorithmique pour prospecter les zones blanches très rentables. Et surtout, avec la Gap Analysis, nous avons créé un cadre quantitatif pour identifier les sous-performances managériales séparément des faiblesses géographiques. Le tout réuni dans un outil qui favorise une prise de décision rapide et visuelle.

%=============================================================================
\section{Slide 38 --- Perspectives}
\subsection{Durée : $\sim$35 secondes}
%=============================================================================
Ce travail ouvre la voie à de fortes évolutions. 
D'abord, diversifier l'enrichissement en intégrant nativement les API d'OpenStreetMap ou Overture pour gagner en indépendance vis-à-vis de Google. Ensuite, la précision des modèles peut être décuplée en passant d'une densité de population à l'échelle de la ville vers des données socio-démographiques à l'échelle du quartier. Enfin, l'objectif ultime serait le déploiement sur l'infrastructure interne de la banque, avec une intégration en temps réel des flux transactionnels pour un ré-entraînement automatique des modèles.

%=============================================================================
\section{Slide 39 --- Merci}
\subsection{Durée : $\sim$10 secondes}
%=============================================================================
J'arrive au terme de ma présentation. Je vous remercie vivement pour votre attention et votre écoute. Je me tiens désormais à votre entière disposition pour répondre à vos questions.

\newpage
%=============================================================================
\begin{center}
{\Large\bfseries\color{accent} Récapitulatif du Timing}
\end{center}
\vspace{0.3cm}

\begin{center}
\renewcommand{\arraystretch}{1.3}
\begin{tabular}{|l|c|c|}
\hline
\textbf{Slide} & \textbf{Durée Estimée} & \textbf{Cumul} \\
\hline
1 à 8 : Introduction, AWB, Objectifs & $\sim$ 2min 45s & 2:45 \\
9 à 14 : Workflow, Data, Technologies & $\sim$ 2min 30s & 5:15 \\
15 à 21 : ML, Régression, Profils & $\sim$ 3min 45s & 9:00 \\
22 à 24 : Sélection Nouveaux Sites & $\sim$ 1min 20s & 10:20 \\
25 à 29 : BCG, Gap Analysis, Fusions & $\sim$ 3min 15s & 13:35 \\
30 à 35 : Démo Dashboard & $\sim$ 1min 35s & 15:10 \\
36 à 39 : Conclusion, Perspectives & $\sim$ 1min 30s & 16:40 \\
\hline
\textbf{Temps de parole pure} & \textbf{$\sim$ 16min 40s} & \\
\textbf{+ Transitions et pauses (marge)} & \textbf{$\sim$ 1min 00s} & \\
\hline
\textbf{TOTAL ESTIMÉ EN SOUTENANCE} & \textbf{$\sim$ 17min 30s} & \\
\hline
\end{tabular}
\end{center}

\vspace{0.3cm}
{\color{tip}\textbf{Note Stratégique :}}
\begin{itemize}
  \item Le temps de parole pure (lecture du script à un rythme posé et naturel) prendra désormais un peu plus de \textbf{16 minutes} sans problème. 
  \item Si vous avez tendance à parler plus vite avec le stress, cette densité de texte vous assure de ne pas tomber sous la barre des 15 minutes. Vous pouvez donc parler sereinement, les détails ajoutés apporteront une vraie crédibilité technique face au jury.
\end{itemize}

\end{document}
